\documentclass[10pt,a4paper]{article}
\usepackage[T1]{fontenc}
\usepackage{lmodern}
\usepackage{amsmath,amssymb,amsthm,mathtools}
\usepackage[margin=24mm]{geometry}
\usepackage{microtype}
\usepackage{booktabs,array,longtable}
\usepackage{xcolor}
\usepackage{hyperref}
\usepackage{fancyhdr}
\hypersetup{colorlinks=true,linkcolor=black,citecolor=black,urlcolor=blue,pdfauthor={Matthew J. Colbrook}}
\pagestyle{fancy}
\fancyhf{}
\fancyhead[L]{\small Resolution / MI-22}
\fancyhead[R]{\small 11 September 2026}
\fancyfoot[C]{\thepage}
\setlength{\headheight}{14pt}
\setlength{\parskip}{4pt}
\setlength{\parindent}{0pt}
\newtheorem{theorem}{Theorem}[section]
\newtheorem{proposition}[theorem]{Proposition}
\newtheorem{lemma}[theorem]{Lemma}
\newtheorem{corollary}[theorem]{Corollary}
\theoremstyle{remark}
\newtheorem{remark}[theorem]{Remark}
\newcommand{\M}{\mathbb M}
\newcommand{\C}{\mathbb C}
\newcommand{\R}{\mathbb R}
\newcommand{\tr}{\operatorname{tr}}
\newcommand{\diag}{\operatorname{diag}}
\newcommand{\spec}{\operatorname{spec}}
\newcommand{\rank}{\operatorname{rank}}
\newcommand{\abs}[1]{\lvert #1\rvert}
\newcommand{\norm}[1]{\lVert #1\rVert}
\newcommand{\opnorm}[1]{\lVert #1\rVert_{\infty}}
\newcommand{\schatten}[2]{\lVert #1\rVert_{#2}}
\newcommand{\symmod}{\mathcal S}
\newcommand{\maxmod}{\mathcal M}
\newcommand{\draftnotice}{\begin{center}\small\textit{Proposed mathematical resolution. Complete proof supplied; not submitted or peer reviewed.}\end{center}}

\usepackage{xurl}
\setlength{\emergencystretch}{3em}
\allowdisplaybreaks[1]

\makeatletter
\renewcommand{\@maketitle}{\newpage\null\begin{center}
{\Large\@title\par}\vskip .8em
{\normalsize\begin{tabular}[t]{c}\@author\end{tabular}\par}\vskip .5em
{\small\@date}\end{center}\par\vskip .5em}
\makeatother

\title{MI-22: Resolution}
\author{Matthew J. Colbrook\\
\small Department of Applied Mathematics and Theoretical Physics\\
\small University of Cambridge, Cambridge, United Kingdom\\
\small\href{mailto:m.colbrook@damtp.cam.ac.uk}{m.colbrook@damtp.cam.ac.uk}}
\date{11 September 2026}
\begin{document}
\maketitle
\textbf{Repository status: Solved.}
The complete argument received an independent Codex-agent PASS review
against the exact catalog target.
The original draft was AI-assisted; the reviewed proof below is unchanged.
This is independent agent verification, not external human peer review or
formal proof certification. \href{https://github.com/ajt60gaibb/OpenProblemsInNLA/blob/main/references/colbrook-matrix-2026-09-11/verification/reviews/MI-22-review.md}{Detailed review and scope}.
\par\medskip
% BEGIN REVIEWED PROOF
\section{MI-22: a certified operator-norm counterexample to the weighted log-majorization}
\label{sec:mi22}
The conjecture is
\begin{equation}\label{eq:mi22-target}
 s\bigl(A^t(A\#_t B)B^{1-t}\bigr)\prec_{\log}s(AB),
 \qquad A,B>0,\quad 0\leq t\leq1,
\end{equation}
where singular values are decreasingly ordered and $\#_t$ is the weighted geometric mean. It is Conjecture 1.1 of \cite{GAMA2021}. We disprove even the largest-singular-value inequality. The proof uses explicit rational certificates for two matrix roots; every certificate check is an exact rational comparison.

\begin{theorem}\label{thm:mi22}
A counterexample to \eqref{eq:mi22-target} is
\[
 A=\diag\left(256,\frac1{256},1\right),\qquad
 B=\begin{pmatrix}17&-4&0\\-4&16385&-8192\\0&-8192&4096\end{pmatrix},
 \qquad t=\frac18.
\]
In fact,
\[
 \opnorm{A^{1/8}(A\#_{1/8}B)B^{7/8}}>10900,
 \qquad \opnorm{AB}<10200.
\]
\end{theorem}

We first give the elementary error estimate used to turn rational root approximations into a proof.
\begin{lemma}\label{lem:root-lipschitz}
If $P,Q\geq mI>0$ and $0<\alpha<1$, then
\[
 \opnorm{P^\alpha-Q^\alpha}
 \leq \alpha m^{\alpha-1}\opnorm{P-Q}.
\]
\end{lemma}
\begin{proof}
The scalar integral representation for $x^\alpha$, applied by spectral calculus, and the resolvent identity give
\[
 P^\alpha-Q^\alpha
 =\frac{\sin(\pi\alpha)}{\pi}\int_0^\infty
 u^\alpha(uI+Q)^{-1}(P-Q)(uI+P)^{-1}\,du.
\]
The integral converges in norm. The norm of the integrand is at most
$u^\alpha(u+m)^{-2}\opnorm{P-Q}$. Differentiating the scalar integral formula, or evaluating this positive integral directly, gives
\[
 \frac{\sin(\pi\alpha)}{\pi}\int_0^\infty\frac{u^\alpha}{(u+m)^2}\,du
 =\alpha m^{\alpha-1},
\]
proving the estimate.
\end{proof}

\begin{proof}[Proof of Theorem \ref{thm:mi22}]
The matrix $A$ is positive definite. The leading principal minors of $B$ are $17$, $278529$, and $4096$, so $B>0$ as well. Define the rational matrix
\[
 C=A^{-1/2}BA^{-1/2}
 =\begin{pmatrix}
 17/256&-4&0\\-4&4194560&-131072\\0&-131072&4096
 \end{pmatrix}.
\]
Set $X=C^{1/8}$ and $Y=B^{1/8}$. The left-hand matrix in the theorem is
\[
 L=A^{5/8}XA^{1/2}Y^7,
 \qquad A^{5/8}=\diag(32,1/32,1),\quad A^{1/2}=\diag(16,1/16,1).
\]
Here are rational approximations $R$ to $X$ and $S$ to $Y$, with a common denominator $10^{15}$:
\begingroup\small
\[
 R=10^{-15}\begin{pmatrix}
 563431071954661&-4774979464975&-152620714401404\\
 -4774979464975&6722248446399974&-185449303604146\\
 -152620714401404&-185449303604146&793308473748417
 \end{pmatrix},
\]
\[
 S=10^{-15}\begin{pmatrix}
 1417565567728212&-40084698659651&-79375309604622\\
 -40084698659651&2883518412755210&-1150490885652445\\
 -79375309604622&-1150490885652445&1157680400969996
 \end{pmatrix}.
\]
\endgroup
These displayed rational matrices obey the following exact certificate inequalities:
\begin{align}
 B,C&>2^{-10}I,\label{eq:mi22-cert1}\\
 \frac9{20}I<R<8I,&\qquad \frac9{20}I<S<4I,\label{eq:mi22-cert2}\\
 \norm{R^8-C}_F^2&<10^{-16},\qquad \norm{S^8-B}_F^2<10^{-16}.\label{eq:mi22-cert3}
\end{align}
All positive-definite comparisons in \eqref{eq:mi22-cert1}--\eqref{eq:mi22-cert2} are verified by positive leading principal minors. The two comparisons in \eqref{eq:mi22-cert3} are sums of nine rational squares. The accompanying certificate verifier performs precisely these operations, without floating-point arithmetic, and its full input is printed above.

Since $(9/20)^8>2^{-10}$, both $R^8$ and $S^8$ exceed $2^{-10}I$. Apply Lemma \ref{lem:root-lipschitz} with $\alpha=1/8$ and $m=2^{-10}$. Because $\alpha m^{\alpha-1}=2^{23/4}<64$, equations \eqref{eq:mi22-cert1}--\eqref{eq:mi22-cert3} imply
\[
 \opnorm{X-R}<\varepsilon,\qquad \opnorm{Y-S}<\varepsilon,
 \qquad \varepsilon=64\cdot10^{-8}.
\]
Also $\opnorm{X}<8$ and $\opnorm{Y}<4$: for example, the exact traces of $C$ and $B$ are respectively smaller than $8^8$ and $4^8$. Together with \eqref{eq:mi22-cert2}, the telescoping identity for $Y^7-S^7$ gives
\[
 \opnorm{Y^7-S^7}\leq7\cdot4^6\varepsilon.
\]
For the rational matrix $\widetilde L=A^{5/8}RA^{1/2}S^7$, it follows that
\begin{align*}
 \opnorm{L-\widetilde L}
 &\leq32\cdot16\,\varepsilon\bigl(4^7+8\cdot7\cdot4^6\bigr)\\
 &=\frac{6291456}{78125}<100.
\end{align*}
One further exact rational multiplication yields
\begin{equation}\label{eq:mi22-entry}
 \widetilde L_{12}>11000.
\end{equation}
Consequently $\opnorm{L}\geq|L_{12}|>11000-100=10900$.

On the other hand,
\[
 AB=\begin{pmatrix}4352&-1024&0\\-1/64&16385/256&-32\\0&-8192&4096\end{pmatrix},
\]
and its exact squared Frobenius norm is
\[
 \norm{AB}_F^2=\frac{6807858741265}{65536}<10200^2.
\]
Thus $\opnorm{AB}\leq\norm{AB}_F<10200<10900<\opnorm{L}$, contradicting the first inequality required by log-majorization.
\end{proof}

\begin{remark}[Reproducibility and the role of computation]
High-precision calculations were used to propose the rational matrices $R,S$, but no such calculation is trusted by the proof. The finite certificate consists of \eqref{eq:mi22-cert1}--\eqref{eq:mi22-cert3}, \eqref{eq:mi22-entry}, and the Frobenius-norm comparison; all are verified with integer fractions. The error estimate bridges these exact finite checks to the true matrix roots. This is a rigorous computer-assisted counterexample, not a claim based only on a numerical singular-value comparison.
\end{remark}

% END REVIEWED PROOF
\begin{thebibliography}{1}

\bibitem{GAMA2021}
Mohammad~M. Ghabries, Hassane Abbas, Bassam Mourad, and Abdallah Assi.
\newblock New log-majorization results concerning eigenvalues and singular
  values and a complement of a norm inequality.
\newblock arXiv:2105.13356, 2021.
\newblock Conjectures 1.1 and 2.1.
\newblock URL: \url{https://arxiv.org/abs/2105.13356}.

\end{thebibliography}

\end{document}
